\subsection{Gradient Boosted Trees and Rule Ensembles}
\begin{definition} \label{def:sl-additive-tree-expansion} \hlt{Additive Tree Expansion}\\
Let $R_1,\ldots,R_J$ be terminal regions, $\gamma_j\in\R$, and $\Theta=\{(R_j,\gamma_j):1\le j\le J\}$.
\begin{enumerate}[label=(\roman*)]
\setlength{\itemsep}{0pt}
\item \hlt{Tree}: $T(x;\Theta)=\sum\limits_{j=1}^J\gamma_j\mathbf{1}_{\{x\in R_j\}}$.
\item \hlt{Additive expansion}: $f_M(x)=\sum\limits_{m=1}^M T(x;\Theta_m)$.
\item \hlt{Terminal fit}: with regions $R_{jm}$ fixed, $\widehat{\gamma}_{jm}\in\operatorname*{argmin}\limits_{\gamma}\sum\limits_{x_i\in R_{jm}}L\{y_i,f_{m-1}(x_i)+\gamma\}$.
\end{enumerate}
\end{definition}

\begin{breakablealgorithm}
\caption{Gradient Tree Boosting}
\label{alg:sl-gradient-tree-boosting}
\begin{algorithmic}
\Require Training data $(x_i,y_i)$, differentiable loss $L(y,f)$, tree sizes $J_m$, learning rate $0<\nu\leq1$, iterations $M$
\Ensure Additive predictor $f_M$
\State $\displaystyle f_0(x)\gets\operatorname*{arg\,min}_{\gamma}\sum_{i=1}^n L(y_i,\gamma)$
\For{$m=1,\ldots,M$}
\For{$i=1,\ldots,n$}
\State $\displaystyle r_{im}\gets-\left[\frac{\partial L(y_i,f(x_i))}{\partial f(x_i)}\right]_{f=f_{m-1}}$
\EndFor
\State Fit a regression tree to targets $r_{im}$, giving terminal regions $R_{jm}$, $j=1,\ldots,J_m$
\For{$j=1,\ldots,J_m$}
\State $\displaystyle \gamma_{jm}\gets
\operatorname*{arg\,min}_{\gamma}
\sum_{x_i\in R_{jm}}L\{y_i,f_{m-1}(x_i)+\gamma\}$
\EndFor
\State $\displaystyle f_m(x)\gets f_{m-1}(x)+
\nu\sum_{j=1}^{J_m}\gamma_{jm}\mathbf{1}_{\{x\in R_{jm}\}}$
\EndFor
\State \textbf{return} $f_M$
\end{algorithmic}
\end{breakablealgorithm}

\begin{remark} \label{rem:sl-gradient-boosting-losses} \hlt{Pseudo-Residuals}\\
At stage $m$, pseudo-residuals are negative pointwise gradients $r_{im}\defeq-\partial_f L(y_i,f)|_{f=f_{m-1}(x_i)}$.\\
Gradient boosting fits the next tree to $(x_i,r_{im})$.
\begin{enumerate}[label=(\roman*)]
\setlength{\itemsep}{0pt}
\item \hlt{Squared error}: $r_{im}=y_i-f_{m-1}(x_i)$.
\item \hlt{Absolute error}: $r_{im}=\operatorname{sign}\{y_i-f_{m-1}(x_i)\}$, away from zero residuals.
\item \hlt{Multinomial deviance}: $r_{ikm}=\mathbf{1}_{\{y_i=\mathcal G_k\}}-p_k(x_i)$, where $p_k$ is the current softmax probability.
\end{enumerate}
\end{remark}

\begin{breakablealgorithm}
\caption{Gradient Boosting for $K$-Class Classification}
\label{alg:sl-multiclass-gradient-boosting}
\begin{algorithmic}
\Require Encoded labels $y_{ik}\in\{0,1\}$, classes $k=1,\ldots,K$, tree sizes $J_m$, iterations $M$
\Ensure Class-score functions $\widehat{f}_1,\ldots,\widehat{f}_K$
\State Initialize $f_{k0}(x)\gets0$ for $k=1,\ldots,K$
\For{$m=1,\ldots,M$}
\State $\displaystyle p_k(x)\gets\frac{\exp(f_{k,m-1}(x))}{\sum_{\ell=1}^K\exp(f_{\ell,m-1}(x))}$ for $k=1,\ldots,K$
\For{$k=1,\ldots,K$}
\State Compute $r_{ikm}\gets y_{ik}-p_k(x_i)$ for $i=1,\ldots,n$
\State Fit a regression tree to targets $r_{ikm}$, giving terminal regions $R_{jkm}$, $j=1,\ldots,J_m$
\For{$j=1,\ldots,J_m$}
\State $\displaystyle \gamma_{jkm}\gets
\frac{K-1}{K}\,
\frac{\sum_{x_i\in R_{jkm}} r_{ikm}}
{\sum_{x_i\in R_{jkm}} |r_{ikm}|(1-|r_{ikm}|)}$
\EndFor
\State $\displaystyle f_{km}(x)\gets f_{k,m-1}(x)+\sum_{j=1}^{J_m}\gamma_{jkm}\mathbf{1}_{\{x\in R_{jkm}\}}$
\EndFor
\EndFor
\State \textbf{return} $\widehat{f}_k(x)=f_{kM}(x)$ for $k=1,\ldots,K$
\end{algorithmic}
\end{breakablealgorithm}

\begin{remark} \label{rem:sl-boosted-tree-regularization} \hlt{Tree Size, Shrinkage, and Subsampling}\\
Boosted-tree complexity is controlled by iterations $M$, terminal counts $J_m$, shrinkage $\nu$, and sampling.
\begin{enumerate}[label=(\roman*)]
\setlength{\itemsep}{0pt}
\item \hlt{Iterations}: larger $M$ adds more trees and increases model capacity.
\item \hlt{Tree size}: a common $J$ caps interactions; stumps ($J=2$) give main effects, and larger $J$ allows order up to $J-1$.
\item \hlt{Shrinkage}: use $f_m=f_{m-1}+\nu T_m$ with $0<\nu\le1$; smaller $\nu$ usually requires larger $M$.
\item \hlt{Subsampling}: fit $T_m$ on random observations to reduce cost and add averaging-like randomness.
\end{enumerate}
These tuning choices trade optimization speed, interaction depth, and test risk.
\end{remark}

\begin{method} \label{meth:sl-boosted-tree-interpretation} \hlt{Relative Importance and Partial Dependence}\\
Let $f$ be a fitted tree ensemble.  For $S\subset\{1,\ldots,p\}$, write $C=S^c$ and $x_i=(x_{iS},x_{iC})$.
\begin{enumerate}[label=(\roman*)]
\setlength{\itemsep}{0pt}
\item \hlt{Relative importance}: for predictor $j$, sum the split-improvement scores over tree nodes splitting on $j$.\\
For an ensemble, average over trees and rescale if desired.
\item \hlt{Partial dependence}: define $f_S(x_S)=\mathbb{E}_{X_C}\{f(x_S,X_C)\}$ and estimate $\widehat f_S(x_S)=n^{-1}\sum\limits_{i=1}^n f(x_S,x_{iC})$.
\end{enumerate}
Partial dependence averages over the marginal distribution of $X_C$.\\
It equals $\mathbb{E}\{f(X)\mid X_S=x_S\}$ only under compatible dependence structure.
\end{method}

\begin{principle} \label{prin:sl-bet-on-sparsity} \hlt{Bet on Sparsity}\\
Let $\mathcal T=\{T_1,\ldots,T_K\}$ and $f_\alpha=\sum\limits_{k=1}^K\alpha_kT_k$, with support $S(\alpha)=\{k:\alpha_k\neq0\}$.\\
An $\ell^1$ penalty is useful when a low-risk predictor has $|S(\alpha)|\ll K$ and $|S(\alpha)|$ is small relative to $n$ and noise.\\
Dense signals with many small $\alpha_k$ remain hard to estimate.\\
Sparsity is meaningful only relative to $n$, noise level, and $\mathcal T$.
\end{principle}

\begin{method} \label{meth:sl-lasso-postprocessed-ensembles} \hlt{Lasso Post-Processing of Ensembles}\\
Let a fitted ensemble generate the fixed dictionary $\mathcal T_L=\{T_1,\ldots,T_M\}$ and let $L$ be the chosen prediction loss.
\begin{enumerate}[label=(\roman*)]
\setlength{\itemsep}{0pt}
\item \hlt{Fit path}: compute $\widehat{\alpha}(\lambda)\in\operatorname*{argmin}\limits_{\alpha}\{\sum\limits_{i=1}^n L(y_i,\alpha_0+\sum\limits_{m=1}^M\alpha_mT_m(x_i))+\lambda\sum\limits_{m=1}^M|\alpha_m|\}$.
\item \hlt{Select}: choose $\lambda$ by validation or an information criterion.
\item \hlt{Return}: use $\widehat f_\lambda(x)=\widehat\alpha_0+\sum\limits_{m:\widehat\alpha_m(\lambda)\neq0}\widehat\alpha_m(\lambda)T_m(x)$.
\end{enumerate}
The post-processed model can use fewer trees.\\
A useful generated dictionary is relevant and not nearly duplicated, with subsampling one way to diversify it.
\end{method}

\begin{figure}[H]
\centering
\begin{tikzpicture}[
    >=Stealth,
    font=\scriptsize,
    line cap=round,
    line join=round,
    box/.style={draw=gray!65, rounded corners=2pt, align=center, inner xsep=6pt, inner ysep=2pt, fill=gray!8},
    generator/.style={box, minimum width=1.35cm},
    dictionary/.style={box, draw=green!45!black, fill=green!8, minimum width=2.45cm},
    selection/.style={box, draw=orange!70!black, fill=orange!10, minimum width=1.8cm},
    predictor/.style={box, minimum width=3.1cm},
    flow/.style={->, draw=black!75, line width=0.65pt}
]
\node[generator] (gen) at (0,0) {$\mathcal A_{\mathrm{ens}}$};
\path (gen.east) ++(0.75,0) node[dictionary, anchor=west] (dict)
{$\mathcal T_L=\{T_1,\ldots,T_M\}$};
\path (dict.east) ++(0.75,0) node[selection, anchor=west] (lasso)
{$\widehat S_\lambda=\{2,7,19\}$};
\path (lasso.east) ++(0.75,0) node[predictor, anchor=west] (pred)
{$\widehat f_\lambda(x)=\sum_{m\in\widehat S_\lambda}\widehat\beta_mT_m(x)$};

\draw[flow] (gen) -- (dict);
\draw[flow] (dict) -- (lasso);
\draw[flow] (lasso) -- (pred);
\end{tikzpicture}

\caption{An ensemble generator $\mathcal A_{\mathrm{ens}}$ builds the \textcolor{green!45!black}{dictionary $\mathcal T_L$}; lasso selects the \textcolor{orange!80!black}{active set $\widehat S_\lambda$} for the sparse predictor $\widehat f_\lambda$.}
\label{fig:sl-ensemble-lasso-postprocessing}
\end{figure}

\begin{method} \label{meth:sl-rule-ensembles} \hlt{Rule Ensembles}\\
Let $T_1,\ldots,T_M$ be fitted trees.  A node path $p$ defines the rule $r_p(x)=\mathbf{1}_{\{x\text{ satisfies all splits on }p\}}$.
\begin{enumerate}[label=(\roman*)]
\setlength{\itemsep}{0pt}
\item \hlt{Extract}: form $\mathcal T^m_{\mathrm{rule}}=\{r_p:p\text{ is a node path in }T_m\}$.
\item \hlt{Aggregate}: set $\mathcal T_{\mathrm{rule}}=\bigcup\limits_{m=1}^M\mathcal T^m_{\mathrm{rule}}$.
\item \hlt{Fit}: regress $y$ on $\mathcal T_{\mathrm{rule}}$ by lasso or another regularized linear method.
\item \hlt{Augment}: optionally add original variables $x_j$ as linear terms.
\end{enumerate}
Rules encode split interactions; added linear terms encode main linear effects.
\end{method}
